\subsection{测度的逆系统}

定义 1. --- 设 \(\mathcal{T} = (\mathrm{T}_i, p_{ij})\) 是以 I 为指标的拓扑空间逆系统。称测度的逆系统（分别地，测度的子逆系统）

为 \(\mathcal{T}\) 上的一个族 \((\mu_i)_{i\in I}\)，其中 \(\mu_i\) 是 \(\mathrm{T}_i\) 上的有界测度，对所有 \(i\in \mathrm{I}\) 均如此，并且 \(\mu_i = p_{ij}(\mu_j)\)（分别地，\(\mu_i\geqslant p_{ij}(\mu_j)\)）对 \(i\leqslant j\) 成立。

命题 3. --- 设给定以 I 为指标的拓扑空间逆系统 \(\mathcal{T} = (\mathrm{T}_i, p_{ij})\)、拓扑空间 T、相容且分离的连续映射族 \(p_i: \mathrm{T} \to \mathrm{T}_i\)（\(i \in \mathrm{I}\)）以及测度子逆系统 \((\mu_i)_{i \in \mathrm{I}}\)（它是在 \(\mathcal{T}\) 上的）。对 T 的每个紧子集 K，置

\[
\mathrm{J} (\mathrm{K}) = \inf _ {i \in \mathrm{I}} \mu_ {i} ^ {\bullet} (p _ {i} (\mathrm{K})).\tag{1}
\]

则 T 上存在且仅存在一个有界测度 \(\pi\)，使得对 T 的每个紧子集 K 都有 \(\pi^{\bullet}(\mathrm{K}) = \mathrm{J}(\mathrm{K})\)。有 \(\mu_i \geqslant p_i(\pi)\) 对所有 \(i \in I\) 成立，并且 \(\pi\) 是满足此条件的 T 上最大测度。

先证明 \(\mathrm{J}(\mathrm{K})\) 是 \(\mu_i^\bullet (p_i(\mathrm{K}))\) 关于有向预序集 I 的截段滤子 \(\mathfrak{F}\) 的极限：为此，只须（GT, IV, §5, No.~2, Th. 2）证明 \(\mu_i^\bullet (p_i(\mathrm{K})) \geqslant \mu_j^\bullet (p_j(\mathrm{K}))\) 对 \(i \leqslant j\) 成立；现在置 \(\mu_{ij}' = p_{ij}(\mu_j)\)，则有 \(\mu_{ij}' \leqslant \mu_i\) 以及 \(p_j(\mathrm{K}) \subset \overline{p}_{ij}^{-1}(p_i(\mathrm{K}))\)，从而

\[
\mu_ {j} ^ {\bullet} \big (p _ {j} (\mathrm{K}) \big) \leqslant \mu_ {j} ^ {\bullet} \big (\overline{p}_{ij}^{-1} (p _ {i} (\mathrm{K})) \big) = (\mu_ {i j} ^ {\prime}) ^ {\bullet} \big (p _ {i} (\mathrm{K}) \big) \leqslant \mu_ {i} ^ {\bullet} \big (p _ {i} (\mathrm{K}) \big).
\]

现在来研究函数 J 的性质：

\begin{enumerate}
\def\labelenumi{\arabic{enumi})}
\item
  显然，\(\mathrm{J}(\mathbf{K})\leqslant \mathrm{J}(\mathbf{L})\) 当 \(\mathbf{K}\subset \mathbf{L}\) 时成立。
\item
  设 \(\mathbf{K}\) 和 \(\mathbf{L}\) 是 \(\mathbf{T}\) 的两个紧子集。对每个 \(i \in \mathbf{I}\)，有 \(p_i(\mathbf{K} \cup \mathbf{L}) = p_i(\mathbf{K}) \cup p_i(\mathbf{L})\)，从而
\end{enumerate}

\[
\mu_ {i} ^ {\bullet} \left(p _ {i} (\mathrm{K} \cup \mathrm{L})\right) \leqslant \mu_ {i} ^ {\bullet} \left(p _ {i} (\mathrm{K})\right) + \mu_ {i} ^ {\bullet} \left(p _ {i} (\mathrm{L})\right);
\]

关于滤子 \(\mathfrak{F}\) 取极限，得到 \(\mathrm{J}(\mathrm{K} \cup \mathrm{L}) \leqslant \mathrm{J}(\mathrm{K}) + \mathrm{J}(\mathrm{L})\)。

\begin{enumerate}
\def\labelenumi{\arabic{enumi})}
\setcounter{enumi}{2}
\tightlist
\item
  假设紧集 K 和 L 不相交。由第 1 号命题 2，存在 \(i \in I\)，使得 \(p_j(K) \cap p_j(L) = \varnothing\) 对 \(j \geqslant i\) 成立。因此对 \(j \geqslant i\) 有
\end{enumerate}

\[
\mu_ {j} ^ {\bullet} \left(p _ {j} (\mathrm{K} \cup \mathrm{L})\right) = \mu_ {j} ^ {\bullet} \left(p _ {j} (\mathrm{K})\right) + \mu_ {j} ^ {\bullet} \left(p _ {j} (\mathrm{L})\right),
\]

从而关于滤子 F 取极限得到 \(\mathrm{J}(\mathrm{K}\cup\mathrm{L})=\mathrm{J}(\mathrm{K})+\mathrm{J}(\mathrm{L})\)。

\begin{enumerate}
\def\labelenumi{\arabic{enumi})}
\setcounter{enumi}{3}
\tightlist
\item
  设 \((\mathbf{K}_{\alpha})_{\alpha \in \mathbf{A}}\) 是 T 的紧子集的一个递减有向族，其交为 K。由第 1 节第 1 号命题，\(p_i(\mathbf{K}) = \bigcap_{\alpha \in \mathbf{A}} p_i(\mathbf{K}_{\alpha})\)，并且
\end{enumerate}

因此 \(\mu_i^\bullet (p_i(\mathbf{K})) = \inf_{\alpha \in \mathbf{A}}\mu_i^\bullet (p_i(\mathbf{K}_\alpha))\) 对所有 \(i\in I\) 成立（§1, No.~6, Prop. 5 的推论）。由此推出

\[
\begin{array}{r l} \mathrm{J(K)} & = \inf _ {i \in \mathrm{I}} \mu_ {i} ^ {\bullet} (p _ {i} (\mathrm{K})) = \inf _ {i \in \mathrm{I}} \inf _ {\alpha \in \mathrm{A}} \mu_ {i} ^ {\bullet} (p _ {i} (\mathrm{K} _ {\alpha})) \\ & = \inf _ {\alpha \in \mathrm{A}} \inf _ {i \in \mathrm{I}} \mu_ {i} ^ {\bullet} (p _ {i} (\mathrm{K} _ {\alpha})) = \inf _ {\alpha \in \mathrm{A}} \mathrm{J(K} _ {\alpha}). \end{array}
\]

\begin{enumerate}
\def\labelenumi{\arabic{enumi})}
\setcounter{enumi}{4}
\tightlist
\item
  取定 \(i \in I\)，置 \(c = \mu_i^\bullet(T_i)\)。则 \(c\) 有限且 \(J(K) \leqslant \mu_i^\bullet(p_i(K)) \leqslant \mu_i^\bullet(T_i)\)，故有 \(J(K) \leqslant c\)，对每个紧集 \(K\)（\(T\) 中）成立。前述性质允许应用 §3 第 1 号定理；于是存在且仅存在一个有界测度 \(\pi\) 定义在 \(T\) 上，使 \(\pi^\bullet(K) = J(K)\) 对每个紧子集 \(K\)（\(T\) 中）成立。对每个 \(i \in I\)，记 \(\nu_i\) 为 \(T_i\) 上的测度，它是 \(\pi\) 在 \(p_i\) 下的像。取 \(i \in I\)、\(A\) 为 \(T_i\) 的紧子集，并令 \(\mathfrak{L}\) 为 \(\overline{p}_i^1(A)\) 的紧子集所成的集合。由 §1 第 2 号备注 3，有 \(\pi^\bullet(\overline{p}_i^1(A)) = \sup_{K \in \mathfrak{L}} \pi^\bullet(K)\)；此外，\(\nu_i^\bullet(A) = \pi^\bullet(\overline{p}_i^1(A))\)，且有 \(J(K) = \pi^\bullet(K)\) 对 \(K \in \mathfrak{L}\) 成立，从而 \(\nu_i^\bullet(A) = \sup_{K \in \mathfrak{L}} J(K)\)。对 \(K \in \mathfrak{L}\)，有 \(p_i(K) \subset A\)，从而 \(J(K) \leqslant \mu_i^\bullet(p_i(K)) \leqslant \mu_i^\bullet(A)\)，最后得到 \(\nu_i^\bullet(A) \leqslant \mu_i^\bullet(A)\)。由于 \(A\) 是 \(T_i\) 中任意的紧集，故 \(\nu_i \leqslant \mu_i\)。命题的最后一个断言显然成立。
\end{enumerate}

证毕。

定理 1（Prokhorov）。 --- 设 \(\mathcal{T} = (\mathrm{T}_{i}, p_{ij})\) 是以 I 为指标的拓扑空间逆系统，T 是拓扑空间，\((p_{i})_{i \in I}\) 是相容且分离的连续映射族 \(p_{i}: \mathrm{T} \to \mathrm{T}_{i}\)。最后，设 \((\mu_{i})_{i \in I}\) 是 \(\mathcal{T}\) 上的测度逆系统。

为使存在有界测度 \(\mu\) 定义在 \(\mathbf{T}\) 上，满足 \(p_i(\mu) = \mu_i\) 对所有 \(i \in \mathbf{I}\) 成立，以下条件成立是必要且充分的：

\begin{enumerate}
\def\labelenumi{(\Alph{enumi})}
\setcounter{enumi}{15}
\tightlist
\item
  对每个 \(\varepsilon > 0\)，存在紧子集 \(K\)（属于 \(T\)），使得 \(\mu_i^\bullet(T_i - p_i(K)) \leqslant \varepsilon\) 对所有 \(i \in I\) 成立。
\end{enumerate}

若此条件成立，测度 \(\mu\) 唯一确定，并且

\[
\mu^ {\bullet} (\mathbf {K}) = \inf _ {i} \mu_ {i} ^ {\bullet} (p _ {i} (\mathbf {K}))\tag{2}
\]

对每个紧集 \(\mathbf{K}\)（属于 \(\mathbf{T}\)）都成立。

先证明 \(\mu\) 的唯一性。设 \(\mu\) 是 T 上的有界测度，满足 \(p_i(\mu) = \mu_i\) 对所有 \(i \in I\) 成立。设 K 是 T 的紧子集；由第 1 节第 2 号命题，集合 K 是 T 的闭子集递减有向族 \(\left(\overline{p_i^1}(p_i(K))\right)_{i \in I}\) 的交。由 §1 第 6 号命题 5 的推论，因而有

\[
\mu^ {\bullet} (\mathrm{K}) = \inf _ {i \in \mathrm{I}} \mu^ {\bullet} \left(\bar {p} _ {i} ^ {- 1} (p _ {i} (\mathrm{K}))\right) = \inf _ {i \in \mathrm{I}} \mu_ {i} ^ {\bullet} \left(p _ {i} (\mathrm{K})\right),
\]

这就证明了公式 (2)。由于在紧集全体上相同的两个测度相等（§1, No.~2, Prop. 2 的推论），故 \(\mu\) 唯一。

由命题 3，存在有界测度 \(\pi\) 定义在 \(\mathbf{T}\) 上，使得 \(\pi^{\bullet}(\mathbf{K}) = \inf_{i\in I}\mu_i^\bullet (p_i(\mathbf{K}))\) 对每个紧子集 \(\mathbf{K}\)（属于 \(\mathbf{T}\)）成立。由公式 (2)，存在有界测度 \(\mu\) 定义在 \(\mathbf{T}\) 上，满足 \(p_i(\mu) = \mu_i\) 对所有 \(i\in I\) 成立，因而等价于关系：

\[
\left(\mathrm{P} ^ {\prime}\right) p _ {i} (\pi) = \mu_ {i} \text { for all } i \in \mathrm{I}.
\]

当 \(i \leqslant j\) 时，有 \(\mu_i = p_{ij}(\mu_j)\)，从而 \(\mu_i^\bullet (\mathbf{T}_i) = \mu_j^\bullet (\mathbf{T}_j)\)；由于 I 是有向的，存在有限数 \(c \geqslant 0\)，使得 \(\mu_i^\bullet (\mathbf{T}_i) = c\) 对所有 \(i \in I\) 成立。由命题 3，测度 \(\mu_i - p_i(\pi)\) 为正，因此它为零当且仅当其总质量为零，即 \(\mu_i(\mathbf{T}_i) = p_i(\pi)^\bullet (\mathbf{T}_i)\)。由于 \(p_i(\pi)^\bullet (\mathbf{T}_i) = \pi^\bullet (\mathbf{T})\)，条件 \((\mathbf{P}')\) 遂等价于 \(\pi^\bullet (\mathbf{T}) = c\)，即（§1, No.~2, 备注 3）等价于性质：

\((\mathbf{P}^{\prime \prime})\sup_{K\in \mathfrak{K}}\pi^{\bullet}(\mathbf{K}) = c\)，其中 \(\mathfrak{K}\) 是 T 的紧子集所成的集合。

现在，对于 \(\mathbf{K} \in \mathfrak{K}\) ，有

\[
\pi^ {\bullet} (\mathrm{K}) = \inf _ {i \in \mathrm{I}} \mu_ {i} ^ {\bullet} (p _ {i} (\mathrm{K})) = c - \sup _ {i \in \mathrm{I}} \mu_ {i} ^ {\bullet} (\mathrm{T} _ {i} - p _ {i} (\mathrm{K}))
\]

而此公式立即推出 (P) 与 \((\mathbf{P}^{\prime \prime})\) 等价。

Q.E.D.

设 \((\mathrm{T}_{i}, p_{ij})\) 是拓扑空间的逆系统。置 \(\mathrm{T} = \varprojlim \mathrm{T}_{i}\)，并记 \(p_{i}\) 为从 \(\mathrm{T}\) 到 \(\mathrm{T}_{i}\) 的典范映射。推广第 III 章 §4 第 5 号定义 2，称 \(\mu\) 为定义在 \(\mathrm{T}\) 上的有界测度，是测度逆系统 \((\mu_{i})_{i \in \mathrm{I}}\) 的逆极限，如果 \(\mu_{i} = p_{i}(\mu)\) 对所有 \(i \in \mathrm{I}\) 成立。定理 1 给出了逆极限测度存在的判据。当空间 \(\mathrm{T}_{i}\) 紧且映射 \(p_{ij}\) 为满射时，\(\mathrm{T}\) 紧，并且 \(p_{i}(\mathrm{T}) = \mathrm{T}_{i}\) 对每个 \(i \in \mathrm{I}\) 成立；因此条件 (P) 得到满足，此时我们得到第 III 章 §4 第 5 号命题 8(iv)。

备注。 --- 设 \((\mu_{i})_{i\in\mathbb{I}}\) 是空间逆系统 \(\mathcal{T}=(\mathrm{T}_{i},p_{ij})\) 上的测度逆系统。设给定拓扑空间 \(T'\) 及连续映射 \(p_{i}^{\prime}:T^{\prime}\to T_{i}\)；假设族 \((p_{i}^{\prime})_{i\in\mathbb{I}}\) 相容，但不一定分离。如果族 \((p_{i}^{\prime})_{i\in\mathbb{I}}\) 满足 Prokhorov 条件 (P)，则存在测度 \(\mu^{\prime}\)（不一定唯一），定义在 \(T^{\prime}\) 上，使得 \(p_{i}^{\prime}(\mu^{\prime})=\mu_{i}\) 对所有 \(i\in I\) 成立。

为此，置 \(\mathbf{T} = \varprojlim \mathbf{T}_i\) 和 \(p' = (p_i')_{i\in I}\)，并记 \(p_i\) 为从 \(\mathbf{T}\) 到 \(\mathbf{T}_i\) 的典范映射；Prokhorov 条件由 \(\mathbf{T}\) 与 \(p_i\) 满足，因为 \(p_i(p'(K')) = p'_i(K')\)，且 \(p'(K')\) 在 \(\mathbf{T}\) 中紧，对 \(\mathbf{K}'\) 的每个紧子集 \(\mathbf{T}'\) 都如此。由定理 1，存在有界测度 \(\mu\) 定义在 \(\mathbf{T}\) 上，满足 \(p_i(\mu) = \mu_i\) 对所有 \(i\in I\) 成立。设 \(\mathbf{K}'\) 是 \(\mathbf{T}'\) 中的紧集；则 \(\mu^\bullet (p'(K')) = \inf_{i\in I}\mu_i^\bullet (p_i'(K'))\)，

从而

\[
\mu^ {\bullet} \left(\mathrm{T} - p ^ {\prime} \left(\mathrm{K} ^ {\prime}\right)\right) = \sup _ {i \in \mathrm{I}} \mu_ {i} ^ {\bullet} \left(\mathrm{T} _ {i} - p _ {i} ^ {\prime} \left(\mathrm{K} ^ {\prime}\right)\right).
\]

令 \(\varepsilon > 0\)；由于 Prokhorov 条件 (P) 由 \(p_i'\) 满足，所以可以找到紧子集 \(K'\)（属于 \(T'\)），使得 \(\mu^\bullet(T - p'(K')) \leqslant \varepsilon\)。于是 §2 第 4 号命题 8 说明存在有界测度 \(\mu'\)，定义在 \(T'\) 上，满足 \(\mu = p'(\mu')\)，从而 \(\mu_i = p_i(\mu) = p_i(p'(\mu')) = p_i'(\mu')\) 对所有 \(i \in I\) 成立。

